% TEMPLATE-MILC-v3.tex - plantilla maestra UNICA de las guias MILC v3.
%
% La usan las tres familias de guias, cada una con su driver:
%   · Grados 6-11  -> scripts/build-guias-g11.py
%   · Bebras       -> scripts/build-guias-bebras.py
%   · Semillero    -> scripts/build-guias-semillero.py
%
% Antes habia tres copias casi identicas (~400 lineas iguales, 6 distintas):
% cada mejora habia que aplicarla tres veces, y en la practica no se hacia
% (el motor de recursos vivia solo en dos de ellas). Lo que varia entre
% familias esta parametrizado en 6 placeholders que cada driver llena
% (se nombran aqui SIN su sintaxis real: el builder reemplaza por texto plano,
% tambien dentro de los comentarios, y un valor multilinea rompe el preambulo):
%
%   HEADER_IZQ        encabezado izquierdo de pagina
%   HEADER_DER        encabezado derecho de pagina
%   PORTADA_ETIQUETA  rotulo sobre el numero grande (GRADO/MOMENTO/MODULO)
%   PORTADA_SERIE     linea de serie de la portada
%   PORTADA_ID        identificador de la guia en la portada
%   PORTADA_CIRCULO   el circulito de grado (°), o vacio si no aplica
%
% El resto de claves las documenta content/guias/_SCHEMA.md. El script valida
% que no queden placeholders sin reemplazar antes de escribir el archivo final.

\documentclass[11pt,letterpaper]{article}
\usepackage[margin=2.0cm]{geometry}
\usepackage[table]{xcolor}
\usepackage{fontspec}
\usepackage[spanish]{babel}
\usepackage{tikz}
% Librerías TikZ para los diagramas de las guías (bloque `recursos:`):
%  · babel  → babel en español vuelve activos caracteres como «>», y TikZ los usa
%             en sus opciones (>=stealth, ->). Sin esta librería, cualquier
%             diagrama con flechas revienta con «\language@active@arg> has an
%             extra }». Debe ir ANTES de que se dibuje nada.
%  · positioning        → sintaxis «below left=2pt and 3pt».
%  · decorations.pathreplacing → llaves ({brace}) para señalar rangos.
\usetikzlibrary{babel,positioning,decorations.pathreplacing}
\usepackage{graphicx}
% Circuitos: el trabajo de electrónica se hace hoy con micro:bit y sus sensores
% integrados (MakeCode, sin cableado), así que no se necesitan esquemáticos.
% Si algún día entran montajes con protoboard, basta descomentar esta línea:
% los diagramas .tex/.tikz declarados en `recursos:` ya se incrustan solos.
% \usepackage{circuitikz}
\usepackage[most]{tcolorbox}
\usepackage{tabularx}
\usepackage{array}
\usepackage{fancyhdr}
\usepackage{titlesec}
\usepackage[document]{ragged2e}  % bandera a la derecha en todo el cuerpo
\usepackage{enumitem}
\usepackage{microtype}

% Helvetica Neue no trae la flecha U+2192, asi que cada «→» escrito literalmente
% en el YAML se compilaba a NADA, en silencio (462 flechas en 61 guias: rutas del
% tipo «paleta → tipografia → lockup» salian rotas). Mapeamos el caracter a su
% version matematica, que si tiene glifo. Asi el autor puede seguir escribiendo
% «→» comodamente en el YAML.
\usepackage{newunicodechar}
\newunicodechar{→}{\ensuremath{\rightarrow}}
% Mismo problema con los simbolos de marca y vinetas: el «✓» de «una palomita ✓»
% salia como cuadrito vacio en el PDF de todas las guias que lo usaban.
\usepackage{pifont}
\newunicodechar{✓}{\ding{51}}
\newunicodechar{✗}{\ding{55}}
\newunicodechar{✔}{\ding{52}}
\newunicodechar{•}{\textbullet}
\newunicodechar{≥}{\ensuremath{\geq}}
\newunicodechar{≤}{\ensuremath{\leq}}

\setmainfont{Helvetica Neue}
\setsansfont{Helvetica Neue}
\setlength{\parindent}{0pt}
\setlength{\parskip}{6pt}
% Sin particion de palabras: en 6.º-7.º una palabra cortada por guion al final
% de linea («Comuni-cacion») frena la lectura mucho mas que una linea corta.
\hyphenpenalty=10000
\exhyphenpenalty=10000
\pretolerance=2000
\tolerance=3000
\setlength{\emergencystretch}{3em}
\linespread{1.10}
\renewcommand{\arraystretch}{1.25}
\setlist{leftmargin=5mm,itemsep=1mm,topsep=1mm}
\newcolumntype{Y}{>{\RaggedRight\arraybackslash}X}

\definecolor{milcMagenta}{HTML}{E6007E}
\definecolor{milcVino}{HTML}{5A0038}
\definecolor{milcTurquesa}{HTML}{0093A5}
\definecolor{milcVerde}{HTML}{58A923}
\definecolor{milcMostaza}{HTML}{E5B400}
\definecolor{milcOcre}{HTML}{B86600}
\definecolor{milcNegro}{HTML}{241816}
\definecolor{milcGris}{HTML}{F7F7F4}
\definecolor{milcLinea}{HTML}{E8E2DD}
\definecolor{milcGradePrimary}{HTML}{0066FF}
\definecolor{milcGradeDark}{HTML}{003D99}
\definecolor{milcGradeSoft}{HTML}{D6E8FF}
\definecolor{milcGradeAccent}{HTML}{CFE7FF}
\definecolor{milcGradeText}{HTML}{FFFFFF}
\color{milcNegro}

\pagestyle{fancy}
\fancyhf{}
\lhead{\small Tecnología e Informática · Grado 6}
\rhead{\small Guía 9 · MILC v3}
\cfoot{\small\thepage}
\renewcommand{\headrulewidth}{0pt}

\newtcolorbox{titlebox}[1]{
  enhanced,colback=#1,colframe=#1,arc=7mm,boxrule=0pt,
  left=7mm,right=7mm,top=3mm,bottom=3mm,
  before skip=8pt,after skip=8pt,
  fontupper=\sffamily\bfseries\Large\color{white}
}

\newtcolorbox{softbox}[3]{
  enhanced,colback=#2,colframe=#1,borderline west={4pt}{0pt}{#1},
  arc=2mm,boxrule=.4pt,left=4mm,right=4mm,top=3mm,bottom=3mm,
  before skip=6pt,after skip=8pt,title={#3},coltitle=white,
  colbacktitle=#1,fonttitle=\sffamily\bfseries,fontupper=\RaggedRight,
  attach boxed title to top left={xshift=3mm,yshift=-2mm},
  boxed title style={arc=2mm, boxrule=0pt}
}

\newtcolorbox{drawbox}[2]{
  enhanced,colback=white,colframe=#1,arc=4mm,boxrule=1pt,
  left=5mm,right=5mm,top=4mm,bottom=4mm,before skip=8pt,after skip=8pt,
  title={#2},coltitle=white,colbacktitle=#1,fonttitle=\sffamily\bfseries,
  fontupper=\RaggedRight,
  attach boxed title to top left={xshift=4mm,yshift=-2mm},
  boxed title style={arc=3mm, boxrule=0pt}
}

\newtcolorbox{infoband}[1]{
  enhanced,colback=#1!8,colframe=#1!50,arc=2mm,boxrule=.5pt,
  left=4mm,right=4mm,top=2mm,bottom=2mm,before skip=4pt,after skip=4pt,
  fontupper=\small\RaggedRight
}

% Puente narrativo entre secciones (contrato editorial Conéctate).
% Da continuidad: "Acabas de X. Ahora vas a Y porque Z."
% Visualmente: banda izquierda en color del grado, texto en itálica pequeña.
\newtcolorbox{puentebox}{
  enhanced,colback=milcGradeSoft,colframe=milcGradeDark,
  borderline west={3pt}{0pt}{milcGradeDark},
  arc=0mm,boxrule=0pt,left=5mm,right=4mm,top=2.5mm,bottom=2.5mm,
  before skip=4pt,after skip=8pt,
  fontupper=\itshape\small\color{milcNegro}\RaggedRight
}
% La flecha va en modo matematico a proposito: Helvetica Neue NO tiene el glifo
% U+2192, asi que \textrightarrow se compilaba a NADA («Missing character: There
% is no → in font Helvetica Neue») y el puente salia sin flecha. En modo
% matematico la flecha viene de la fuente matematica y si se dibuja.
\newcommand{\puente}[1]{%
  \begin{puentebox}%
    \textcolor{milcGradeDark}{$\rightarrow$}\hspace{2mm}#1%
  \end{puentebox}%
}

\newcommand{\linead}[1]{\noindent\color{milcLinea}\rule{#1}{0.5pt}\color{milcNegro}}
\newcommand{\respuesta}[1]{\par\vspace{2mm}\foreach \n in {1,...,#1}{\noindent\color{milcLinea}\rule{\linewidth}{0.5pt}\par\vspace{5mm}}\color{milcNegro}}
\newcommand{\checkbox}{\textcolor{milcGradeDark}{\fbox{\phantom{x}}}\hspace{5pt}}

\newcommand{\phasecard}[4]{
\begin{tcolorbox}[enhanced,colback=#3,colframe=#2,arc=3mm,boxrule=.5pt,height=3.05cm,valign=center,halign=center,left=1.5mm,right=1.5mm,top=2mm,bottom=2mm]
{\sffamily\bfseries\fontsize{22}{24}\selectfont\textcolor{#2}{#1}}\par
{\sffamily\bfseries\small #4}
\end{tcolorbox}
}

% ── Piezas del rediseno 2026 (legibilidad 6.º-11.º) ───────────────────
% Banda de estacion: reemplaza el titlebox macizo. Titulo a la izquierda,
% dato practico (tiempo/modalidad) a la derecha, en una sola linea.
\newcommand{\milcestacion}[2]{%
\begin{tcolorbox}[enhanced,colback=#1,colframe=#1,arc=2mm,boxrule=0pt,
  left=5mm,right=5mm,top=2.6mm,bottom=2.6mm,before skip=11pt,after skip=7pt]
{\sffamily\bfseries\fontsize{15}{18}\selectfont\color{white}#2}
\end{tcolorbox}}

% Tarjeta de paso numerada: sustituye las tablas «Pilar / Aplicacion», que
% eran formato de consulta docente, no de estudiante. El numero va en un
% circulo sobre el borde para que la secuencia se lea de un vistazo.
\newcommand{\milcpaso}[3]{%
\begin{tcolorbox}[enhanced,colback=white,colframe=milcLinea,arc=1.5mm,boxrule=.9pt,
  left=14mm,right=4mm,top=3mm,bottom=3mm,before skip=4pt,after skip=4pt,
  fontupper=\RaggedRight,
  overlay={\node[anchor=north west,circle,fill=#2,text=white,inner sep=0pt,
    minimum size=7.4mm,font=\sffamily\bfseries\small]
    at ([xshift=3.6mm,yshift=-3.2mm]frame.north west) {#1};}]
#3
\end{tcolorbox}}

% Casilla marcable de tamano util con lapiz (la anterior era \fbox{\phantom{x}},
% de alto variable segun el texto que la seguia).
\newcommand{\milcmarca}{\framebox[3.8mm]{\rule{0pt}{3.4mm}}}

% Fila marcable de lista suelta (checks de la estacion 1).
\newcommand{\milccheckrow}[1]{%
\par\vspace{1.8mm}\noindent
\makebox[6.4mm][l]{\raisebox{-.55mm}{\textcolor{milcNegro}{\milcmarca}}}%
\parbox[t]{\dimexpr\linewidth-6.4mm\relax}{\RaggedRight #1}\par}

% Celda marcable dentro de la rubrica (tabla de autoevaluacion). Misma
% sangria francesa, pero pensada para vivir dentro de una columna X.
\newcommand{\milcopcion}[1]{%
\makebox[5.2mm][l]{\raisebox{-.55mm}{\textcolor{milcNegro}{\milcmarca}}}%
\parbox[t]{\dimexpr\linewidth-5.2mm\relax}{\RaggedRight #1}}

% ── Recursos visuales de la guía ──────────────────────────────────────
% Declarados en el bloque `recursos:` del YAML y emitidos por el builder.
% \guiaFigura   → imagen rasterizada o PDF (foto, carta celeste, montaje).
% \guiaDiagrama → fuente TikZ/circuitikz (.tex) incrustada como vector:
%                 es la vía para circuitos y esquemáticos editables.
% #1 = ruta relativa al .tex · #2 = pie (puede ir vacío)
\newcommand{\guiaFigura}[2]{%
\begin{center}
\includegraphics[width=0.88\linewidth,height=0.38\textheight,keepaspectratio]{#1}\\[1.5mm]
{\footnotesize\itshape\textcolor{milcNegro!75}{#2}}
\end{center}
}

\newcommand{\guiaDiagrama}[2]{%
\begin{center}
\input{#1}\\[1.5mm]
{\footnotesize\itshape\textcolor{milcNegro!75}{#2}}
\end{center}
}

\begin{document}
\thispagestyle{empty}

% ── Portada ───────────────────────────────────────────────────────────
% Antes: un rectangulo de color a pagina completa. Se veia bien en pantalla,
% pero en la fotocopiadora del colegio se comia el toner y salia gris sucio.
% Ahora la banda ocupa el tercio superior y el resto va en blanco.
\begin{tikzpicture}[remember picture,overlay]
  \path[fill=milcGradePrimary]
    (current page.north west) rectangle ([yshift=-10.6cm]current page.north east);

  \node[anchor=north west,text=milcGradeText] at ([xshift=2.0cm,yshift=-1.55cm]current page.north west)
    {\sffamily\bfseries\fontsize{12}{14}\selectfont GRADO};
  \node[anchor=north west,text=milcGradeText,opacity=.85] at ([xshift=2.0cm,yshift=-2.35cm]current page.north west)
    {\sffamily\bfseries\fontsize{8.5}{10}\selectfont SERIE GUÍAS · TECNOLOGÍA E INFORMÁTICA};
  \node[anchor=north east,text=milcGradeText,opacity=.85,align=right] at ([xshift=-2.0cm,yshift=-1.55cm]current page.north east)
    {\sffamily\bfseries\fontsize{9}{12}\selectfont I.E. Sor María Juliana\\Cartago, Valle del Cauca};

  \node[anchor=west,text=milcGradeText] (gradonum) at ([xshift=1.85cm,yshift=-7.1cm]current page.north west)
    {\sffamily\bfseries\fontsize{112}{112}\selectfont 6};
  % El circulito de grado (°) solo aplica a las guias de grado («11°»); en
  % Bebras y Semillero el numero grande es un momento/modulo, no un grado.
  % Se posiciona relativo al nodo `gradonum`, no en coordenadas absolutas.
  \draw[milcGradeText,line width=1.6pt]
    ([xshift=2.0mm,yshift=-4.5mm]gradonum.north east) circle (.15cm);

  \node[anchor=west,text=milcGradeText,text width=12.3cm,align=left] at ([xshift=6.5cm,yshift=-6.6cm]current page.north west)
    {
     {\sffamily\bfseries\fontsize{9}{11}\selectfont\MakeUppercase{Período 1 · Comunicación e identidad digital}}\\[3mm]
     {\sffamily\bfseries\fontsize{21}{25}\selectfont Comunicación\\responsable ---\\los 3 filtros\\antes de\\publicar}\\[3.5mm]
     {\sffamily\bfseries\fontsize{15}{18}\selectfont Guía 9-6-TIC}
    };
\end{tikzpicture}

\vspace*{9.4cm}
\vfill

{\sffamily\bfseries\fontsize{8}{10}\selectfont\textcolor{milcGradeDark}{LO QUE TE LLEVAS HOY}}

\begin{tcolorbox}[enhanced,colback=white,colframe=milcNegro,arc=2mm,boxrule=1.6pt,
  left=5mm,right=5mm,top=4mm,bottom=4mm,before skip=3pt,after skip=12pt,
  fontupper=\RaggedRight\fontsize{11}{15.5}\selectfont]
Un \textbf{decálogo personal del comunicador responsable} en una hoja del cuaderno con \textbf{tus 3 filtros} (¿es verdad? ¿es necesario? ¿es amable?), \textbf{2 técnicas para no estar de acuerdo con respeto}, y \textbf{1 estrategia para defender a alguien que están atacando}. Más \textbf{3 mensajes reales} que tú reescribes en el cuaderno aplicando los filtros (un comentario, un chat, una historia). Firma de cierre: \emph{``Antes de pulsar enviar, paso por mis 3 filtros.''}.
\end{tcolorbox}

\begin{tcolorbox}[enhanced,colback=milcGris,colframe=milcGris,arc=2mm,boxrule=0pt,
  left=5mm,right=5mm,top=3.5mm,bottom=3.5mm,after skip=12pt]
{\sffamily\bfseries\fontsize{8}{10}\selectfont\textcolor{milcNegro!55}{ESTA GUÍA ES DE}}\\[3mm]
\begin{tabularx}{\linewidth}{X p{3.2cm} p{3.2cm}}
\linead{\linewidth} & \linead{3.0cm} & \linead{3.0cm}\\[-1mm]
{\fontsize{8.5}{10}\selectfont\textcolor{milcNegro!55}{Nombre completo}} &
{\fontsize{8.5}{10}\selectfont\textcolor{milcNegro!55}{Grupo}} &
{\fontsize{8.5}{10}\selectfont\textcolor{milcNegro!55}{Fecha}}\\
\end{tabularx}
\end{tcolorbox}

\vfill

\noindent\textcolor{milcLinea}{\rule{\linewidth}{1pt}}\\[2mm]
{\sffamily\bfseries\fontsize{9.5}{12}\selectfont Autor: Dr. Álvaro Cárdenas Orozco}\\[1mm]
{\sffamily\fontsize{8.5}{11}\selectfont\textcolor{milcNegro!55}{Metodología MILC · apertura ancestral · pensamiento computacional · triángulo Dussel–estoicismo–Floridi}}

\newpage

\section*{Apertura: Comunicación responsable --- los 3 filtros antes de publicar}

\begin{softbox}{milcGradeDark}{milcGradeSoft}{Saber ancestral}
\textbf{En la cocina de la abuela había un colador.} Si alguna vez viste a tu abuela hacer agua de panela o jugo, sabe que el colador era pieza clave. Por ahí pasaba el agua después de hervir la panela, y el colador atrapaba las piedritas, el polvo y los pedazos. \textbf{Lo que pasa el colador llega a la mesa; lo que no pasa, se bota.} La abuela no botaba todo: solo lo que no servía. \textbf{En la cabeza de las personas sensatas hay otro colador}: pasa por ahí lo que van a decir, lo que van a publicar, lo que van a reenviar. La idea no es callar: la idea es \textbf{filtrar}. Decir lo que sirve, dejar de lado lo que no. En internet ese colador es más importante que en cualquier otra parte, porque lo que sale ya no vuelve.
\end{softbox}

\begin{softbox}{milcTurquesa}{milcGris}{Saber contemporáneo}
Hasta ahora aprendiste mucho sobre \emph{recibir}: identidad, huella, riesgos, privacidad. Hoy le toca al \emph{emitir}: \textbf{lo que TÚ publicas, escribes, reenvías y comentas}. Comunicar responsablemente no significa callar; significa \textbf{pasar por 3 filtros} antes de publicar: \emph{(1) ¿Es verdad? (2) ¿Es necesario? (3) ¿Es amable?}. Si pasa los 3, está listo. Si falla en uno, lo piensas mejor. También vas a aprender 2 técnicas para \textbf{discrepar sin pelear} (puedes no estar de acuerdo y seguir siendo amigo) y 1 estrategia para \textbf{defender} a alguien que están atacando (ser \emph{upstander}, no \emph{bystander}). Estos saberes son tan útiles en internet como en la mesa con tu familia.
\end{softbox}

\begin{softbox}{milcMostaza}{milcGris}{Pregunta-puente}
\textit{Si todo lo que has escrito en chats este mes fuera leído en voz alta el lunes en la izada de bandera del colegio, ¿\textbf{cómo te sentirías}? ¿Habría cosas que cambiarías? ¿Cuáles?}
\end{softbox}

\begin{softbox}{milcVerde}{milcGris}{Saber y saber hacer}
Al terminar la clase: (1) podrás \textbf{aplicar} los 3 filtros a cualquier mensaje antes de enviarlo; (2) sabrás \textbf{explicar} 2 técnicas para no estar de acuerdo con respeto; (3) podrás \textbf{analizar} cuándo es momento de defender a alguien en línea; (4) habrás \textbf{creado} tu decálogo personal y 3 mensajes reescritos.
\end{softbox}



\small
\begin{tabularx}{\textwidth}{>{\bfseries}p{3.0cm}Y}
\rowcolor{milcGradePrimary!12}Autor & Dr. Álvaro Cárdenas Orozco\\
DBA & Reconoce principios y conceptos propios de la tecnología, relacionando artefactos con su utilización segura (MEN, T\&I 6°)\\
Referentes & Saber ancestral del pregonero · Pensamiento computacional inicial · Cuidado de la palabra\\
Producto final & Un \textbf{decálogo personal del comunicador responsable} en una hoja del cuaderno con \textbf{tus 3 filtros} (¿es verdad? ¿es necesario? ¿es amable?), \textbf{2 técnicas para no estar de acuerdo con respeto}, y \textbf{1 estrategia para defender a alguien que están atacando}. Más \textbf{3 mensajes reales} que tú reescribes en el cuaderno aplicando los filtros (un comentario, un chat, una historia). Firma de cierre: \emph{``Antes de pulsar enviar, paso por mis 3 filtros.''}.\\
\end{tabularx}
\normalsize

\puente{Hoy haces 4 cosas: aplicas los 3 filtros a mensajes reales, aprendes a discrepar sin pelear, decides cuándo defender a alguien, y firmas tu decálogo.}

\milcestacion{milcGradeDark}{Ruta MILC: así vamos a aprender}
\begin{tabularx}{\textwidth}{>{\centering\arraybackslash}X>{\centering\arraybackslash}X>{\centering\arraybackslash}X>{\centering\arraybackslash}X}
\phasecard{1}{milcVerde}{milcGris}{Escucha\\[-1mm]\small Observo, escucho y pregunto.} &
\phasecard{2}{milcTurquesa}{milcGris}{Entender\\[-1mm]\small Sistematizo y ordeno.} &
\phasecard{3}{milcMagenta}{milcGris}{Hacer\\[-1mm]\small Construyo y aplico.} &
\phasecard{4}{milcVino}{milcGris}{Cierre\\[-1mm]\small Evalúo y reflexiono.}
\end{tabularx}


\puente{Empezamos analizando 5 mensajes que parecen normales pero rompen los filtros.}

\milcestacion{milcVerde}{Estación 1 de 3 · Escucha}

\begin{softbox}{milcVerde}{milcGris}{Escena para escuchar}
\textbf{Actividad 1 · ANALIZA --- 5 mensajes que parecen normales} (12 min · individual). Lee los 5 mensajes y para cada uno, en el cuaderno, escribe: \textbf{(a) ¿Pasaría el filtro de ``es verdad''?} \textbf{(b) ¿Pasaría el de ``es necesario''?} \textbf{(c) ¿Pasaría el de ``es amable''?}. Después decide \emph{sí, lo envío} o \emph{no, lo reescribo}. \textbf{M1:} \emph{``Vi que Pedro lloró en clase, mándeselo al grupo''}. \textbf{M2:} \emph{``Profe, su clase de hoy estuvo aburridísima.''} (en el chat del salón). \textbf{M3:} \emph{``Reenvío: en Cartago hay un virus que mata por mensajes, no abran links''}. \textbf{M4:} \emph{``Lara, no entendí el ejercicio 3, ¿me lo explicas?''}. \textbf{M5:} \emph{``El uniforme de Sara hoy estaba feo''} (en una historia pública).
\end{softbox}

{\sffamily\bfseries\fontsize{8}{10}\selectfont\textcolor{milcNegro!55}{MARCA CUANDO LO HAYAS HECHO}}
\milccheckrow{Leí los 5 mensajes con calma.}
\milccheckrow{Apliqué los 3 filtros a cada uno.}
\milccheckrow{Decidí cuáles enviaría y cuáles reescribiría, con razones.}

\begin{infoband}{milcVerde}


\textbf{Lo que va al cuaderno (Actividad 1 · ANALIZA).} Título: «Actividad 1 --- 5 mensajes por el colador». \textbf{Formato:} tabla de 5 filas con 4 columnas (mensaje, ¿verdad?, ¿necesario?, ¿amable?). \textbf{Extensión:} 5 filas. \textbf{Sabes que terminaste cuando} tienes claro qué falla en cada mensaje que no pasaría.
\end{infoband}

\puente{Después de los mensajes, aprendes las 3 técnicas: filtrar, discrepar con respeto, ser \emph{upstander}.}

\milcestacion{milcTurquesa}{Estación 2 de 3 · Entender}

\begin{softbox}{milcTurquesa}{milcGris}{Lo que vas a entender}
Los 3 filtros son la base. Pero también necesitas \textbf{2 técnicas más} para los momentos difíciles: cuando no estás de acuerdo con alguien, y cuando ves que están atacando a otro. Hoy aprendes las 3 herramientas.
\end{softbox}

\milcpaso{1}{milcTurquesa}{\textbf{Técnica 1 --- Los 3 filtros antes de publicar.} \emph{Filtro 1: ¿Es verdad?} Verifico que lo que digo o reenvío sea cierto. Si no estoy seguro, no lo mando. Las cadenas falsas, las noticias raras, los rumores --- no pasan. \emph{Filtro 2: ¿Es necesario?} ¿Aporta algo? ¿Suma? Si solo es ruido o chisme, no pasa. \emph{Filtro 3: ¿Es amable?} ¿Lo diría en persona, frente a la persona? ¿Lo diría delante de mi mamá? Si no, no pasa.}
\milcpaso{2}{milcTurquesa}{\textbf{Técnica 2 --- Discrepar sin pelear.} Puedes pensar diferente sin convertir todo en pelea. \emph{Paso A:} antes de responder con rabia, espera 5 minutos. \emph{Paso B:} usa la fórmula \textbf{``Yo siento... cuando... porque...''} en vez de \textbf{``Tú eres... siempre...''}. Ejemplo: en vez de \emph{``Eres un mentiroso''}, di \emph{``Yo siento que no me crees cuando dices eso, porque vi otra cosa''}. \emph{Paso C:} acepta que el otro puede no cambiar de opinión --- y aun así pueden seguir hablando.}
\milcpaso{3}{milcTurquesa}{\textbf{Técnica 3 --- Ser \emph{upstander} (no \emph{bystander}).} Cuando ves que están atacando a alguien en internet, hay 3 opciones: (a) callar y mirar (\emph{bystander} = espectador), (b) unirse al ataque, (c) defender al atacado (\emph{upstander} = el que se para). \textbf{Defender no es pelear:} es escribirle \emph{en privado} a la víctima (\emph{``vi lo que pasa, ¿estás bien?''}), avisar a un adulto si es grave, o decir en el grupo algo neutral como \emph{``Eso no me parece justo''}. Cada vez que tú defiendes, los otros también se animan a hacerlo.}
\milcpaso{4}{milcTurquesa}{\textbf{4 errores en comunicación digital que destruyen amistades.} (1) \textbf{Escribir con rabia caliente} --- siempre se sale peor. La regla es: si tienes rabia, espera. (2) \textbf{Sacar conversaciones de contexto} --- mandar capturas a otros mostrando solo una parte. La parte completa cambia todo. (3) \textbf{Reenviar lo que te contaron en privado} --- rompe la confianza más rápido que nada. (4) \textbf{Sumarte a un ataque por miedo} --- pasa todo el tiempo: alguien empieza a burlarse de un compañero y otros se suman para no quedar mal. Eso te hace cómplice.}

\begin{drawbox}{milcTurquesa}{Las 3 técnicas del comunicador responsable}
\textbf{1.} Los 3 filtros antes de publicar: ¿verdad? + ¿necesario? + ¿amable? \\[1mm]
\textbf{2.} Discrepar sin pelear: \emph{``Yo siento... cuando... porque...''} \\[1mm]
\textbf{3.} Ser upstander, no bystander: cuando atacan a otro, no callas.
\end{drawbox}

\begin{softbox}{milcMostaza}{milcGris}{Errores comunes y su corrección}
\textbf{Frases tóxicas que se repiten en chats y qué decir en cambio.} (1) \emph{``Eres un tonto''} → \emph{``No estoy de acuerdo con lo que hiciste''}. (2) \emph{``Nadie te quiere''} → \emph{``Yo me siento incómodo cuando me dices eso''}. (3) \emph{``Es solo broma''} (cuando otro se queja) → \emph{``Entiendo que te molestó, te pido disculpas''}. (4) \emph{``Si fueras mi amigo me apoyarías''} → \emph{``Yo necesito apoyo, ¿puedes escucharme un momento?''}. La diferencia: \textbf{las frases de la derecha no destruyen relaciones, las cuidan}.
\end{softbox}

\begin{infoband}{milcTurquesa}


\textbf{Lo que va al cuaderno (Actividad 2 · EVALÚA).} Título: «Actividad 2 --- Las 3 técnicas y cuándo cada una funciona mejor». \textbf{Formato:} 3 bloques numerados. Cada uno: nombre de la técnica + cuándo se usa + ejemplo personal + \textbf{cuándo NO funciona}. \textbf{Veredicto final (2 líneas):} en una pelea típica de WhatsApp del salón, ¿\emph{cuál de las 3 técnicas} aplicarías primero y por qué? Justifica con 2 razones. \textbf{Extensión:} 3 párrafos + veredicto. \textbf{Sabes que terminaste cuando} podrías recomendarle a un primo más pequeño cuál técnica usar para su problema concreto.
\end{infoband}

\puente{Con las técnicas en mente, escribes tu decálogo y reescribes 3 mensajes propios.}

\milcestacion{milcMagenta}{Estación 3 de 3 · Hacer}

\begin{softbox}{milcMagenta}{milcGris}{Cómo vas a construir tu producto}
\textbf{Actividad 3 · CREA --- Tu decálogo personal + 3 mensajes reescritos} (28 min · individual). Vas a hacer en una hoja del cuaderno tu \textbf{decálogo del comunicador responsable} (10 reglas tuyas) y vas a \textbf{reescribir 3 mensajes} reales que tú o un compañero hayan enviado, aplicando los 3 filtros y las 2 técnicas.
\end{softbox}

\milcpaso{1}{milcMagenta}{\textbf{Paso 1 --- Encabezado y formato.} En la parte superior escribe el título: \textbf{``Mi decálogo del comunicador responsable''} y debajo: \emph{``[Tu nombre], grado 6[tu letra], año 2026''}. Subraya el título. Numera del 1 al 10 en la hoja antes de empezar a escribir el contenido.}
\milcpaso{2}{milcMagenta}{\textbf{Paso 2 --- Las 10 reglas.} Escribe 10 reglas \textbf{con tus palabras}. Las 3 primeras son los 3 filtros (verdad, necesario, amable). Las siguientes 7 son tuyas: pueden incluir cosas como \emph{``Espero 5 minutos antes de responder con rabia''}, \emph{``No reenvío cadenas''}, \emph{``Si veo a alguien siendo atacado, escribo en privado a la víctima''}, etc. No copies; piensa cuáles aplican a tu vida.}
\milcpaso{3}{milcMagenta}{\textbf{Paso 3 --- Bloque de mensajes reescritos.} Debajo del decálogo, escribe el subtítulo \textbf{``3 mensajes reescritos''}. Vas a tomar 3 mensajes que \textbf{tú o un compañero hayan enviado} en los últimos días (sin nombrarlos), uno por cada tipo: un \emph{comentario en redes}, un \emph{chat de grupo}, una \emph{historia o publicación}.}
\milcpaso{4}{milcMagenta}{\textbf{Paso 4 --- Mensaje 1.} Escribe el mensaje original (sin nombres) en \emph{itálicas}. Debajo escribe: \emph{``Falla el filtro de [verdad/necesario/amable] porque [razón]''}. Después escribe \textbf{tu versión nueva} que sí pasa los 3 filtros. Comparas las 2 versiones uno al lado del otro si quieres.}
\milcpaso{5}{milcMagenta}{\textbf{Paso 5 --- Mensaje 2 y 3.} Repite el Paso 4 con los otros 2 mensajes. Procura que los 3 sean de \emph{distintos contextos}: chat, comentario público, historia. Eso te entrena para los 3 espacios distintos.}
\milcpaso{6}{milcMagenta}{\textbf{Paso 6 --- Firma de cierre.} En la esquina inferior derecha escribe la frase: \textbf{``Antes de pulsar enviar, paso por mis 3 filtros.''}. Subráyala. Firma con nombre completo y fecha. La idea es que veas esa hoja antes de mandar algo que dudes.}

\begin{drawbox}{milcMagenta}{Antes de cerrar tu decálogo}
\checkbox El decálogo tiene 10 reglas numeradas. \\[1mm]
\checkbox Las 3 primeras son los 3 filtros (verdad, necesario, amable). \\[1mm]
\checkbox Las otras 7 son escritas con mis palabras. \\[1mm]
\checkbox Hay 3 mensajes reescritos (uno por cada contexto). \\[1mm]
\checkbox Cada mensaje tiene la versión original y la versión nueva. \\[1mm]
\checkbox La firma de cierre está subrayada con nombre y fecha.
\end{drawbox}

\begin{softbox}{milcMagenta}{milcGris}{Plantilla de trabajo}
\textbf{Estructura del decálogo.} \textit{Título.} \textit{10 reglas numeradas.} \textit{Subtítulo: ``3 mensajes reescritos''.} \textit{3 bloques: mensaje original + filtro que falla + versión nueva.} \textit{Esquina: frase de cierre subrayada + firma.}
\end{softbox}

\begin{infoband}{milcMagenta}


\textbf{Lo que va al cuaderno (Actividad 3 · CREA).} Título: «Actividad 3 --- Mi decálogo del comunicador responsable». \textbf{Formato:} hoja completa con decálogo + 3 mensajes reescritos + firma. \textbf{Extensión:} 1 hoja entera. \textbf{Sabes que terminaste cuando} podrías mostrarle el decálogo a tu mamá y ella diría: \emph{``ese eres tú escribiendo con cabeza''}.
\end{infoband}

\puente{El producto es el decálogo + los 3 mensajes reescritos.}

\milcestacion{milcMostaza}{Producto de la sesión}
\textbf{Decálogo personal + 3 mensajes reescritos · firmado}.
\begin{softbox}{milcMostaza}{milcGris}{Criterios de calidad}
(1) El decálogo tiene exactamente 10 reglas numeradas. (2) Las 3 primeras son los 3 filtros, las 7 siguientes son personalizadas. (3) Hay 3 mensajes reescritos en 3 contextos distintos. (4) Cada uno muestra la versión original y la versión nueva. (5) Está firmado, fechado y con la frase de cierre visible.
\end{softbox}

\puente{Antes de salir, miras tus 3 mensajes nuevos: ¿pasan los 3 filtros sin trampa?}

\milcestacion{milcVino}{Evaluación liberadora · cinco dimensiones}

\begin{softbox}{milcVino}{milcGris}{Sentido de la evaluación}
Evaluar no es solo revisar una nota. Es reconocer qué aprendí, qué cambió en mí, qué emociones manejé, cómo me relaciono con la ciudad y la comunidad, y qué le dejaré a quien viene detrás.
\end{softbox}

\textbf{Mi reflexión liberadora en cinco dimensiones.} Completa cada frase iniciada:

\small
\begin{tabularx}{\textwidth}{>{\bfseries\RaggedRight\arraybackslash}p{3.4cm}Y>{\RaggedRight\arraybackslash}p{4.6cm}}
\rowcolor{milcVino}\color{white}Dimensión & \color{white}\bfseries Mi respuesta (frame starter) & \color{white}\bfseries Algo que descubrí\\
1. Personal & Lo que más cambió en mí fue \linead{6.5cm} & \linead{4.2cm}\\
2. Emocional & La emoción más fuerte fue \linead{2.5cm} y la regulé \linead{3cm} & \linead{4.2cm}\\
3. Ciudadana & En democracia esto importa porque \linead{6cm} & \linead{4.2cm}\\
4. Local (Cartago) & En mi barrio sería útil para \linead{6.5cm} & \linead{4.2cm}\\
5. Intergeneracional & A mi abuela/abuelo le diría \linead{6.5cm} & \linead{4.2cm}\\
\end{tabularx}
\normalsize

\begin{infoband}{milcVino}
\textbf{Para la próxima sustentación.} Vuelve a esta tabla en seis meses. Verás que las respuestas cambian — no porque lo aprendido cambie, sino porque tú cambias. La evaluación liberadora es una conversación contigo a través del tiempo.
\end{infoband}


\puente{Cerramos con 3 voces que te ayudan a entender por qué hablar bien es un acto de valentía.}

\milcestacion{milcGradeDark}{Triángulo de pensamiento · cierre filosófico}

\begin{softbox}{milcVino}{milcGris}{Dussel · lente del nosotros}
\textit{``Decir la verdad sin amabilidad es violencia. Decirla con amabilidad es justicia.''}

Hay personas que se creen \emph{honestas} solo por decir cualquier cosa cruda: \emph{``yo digo lo que pienso''}. Pero decir la verdad sin cuidado no es honestidad, es violencia. \textbf{La verdad amable existe y es posible}: la usan los buenos profes, los médicos sensatos, los amigos verdaderos. La técnica de los 3 filtros te entrena para eso: poder decir lo que piensas sin destruir al otro.

\textbf{¿Estoy confundiendo \emph{``decir lo que pienso''} con violencia? ¿Cómo puedo decir lo mismo con cuidado?}
\end{softbox}
\linead{\linewidth}\\[1mm]
\linead{\linewidth}

\begin{softbox}{milcMostaza}{milcGris}{Estoicismo · Marco Aurelio · cuidado interior}
\textit{``El que habla pensando, manda dos veces. El que habla sin pensar, se arrepiente dos veces.''}

Pensar antes de hablar es más rápido que arrepentirse después. \textbf{Cada minuto pensando ahorra horas explicando.} El comunicador responsable no es lento: es \emph{eficiente}. Filtra de una, dice lo que sirve, no tiene que pedir disculpas todo el tiempo. Tu yo del futuro te lo agradece.

\textbf{¿Estoy ahorrándome arrepentimientos al pensar antes, o estoy gastando energía borrando lo que dije sin pensar?}
\end{softbox}
\linead{\linewidth}\\[1mm]
\linead{\linewidth}

\begin{softbox}{milcTurquesa}{milcGris}{Floridi · lente de la infoesfera}
\textit{``Cada mensaje que envías construye o destruye relaciones, sin punto medio. No hay mensajes neutros.''}

Hay gente que cree que un mensaje es solo un mensaje, que ``no es para tanto''. \textbf{En realidad cada mensaje suma o resta} a una relación. Los chats con tus amigos del salón están hechos de \emph{cientos de mensajes pequeños}, cada uno construyó o desgastó confianza. Si vas mensaje por mensaje cuidando, la amistad crece. Si vas descuidando, se va apagando sin que nadie sepa cuándo.

\textbf{Mis últimos 10 mensajes en mi grupo principal, ¿\textbf{sumaron o restaron} a las relaciones de ese grupo?}
\end{softbox}
\linead{\linewidth}\\[1mm]
\linead{\linewidth}

\begin{infoband}{milcNegro}
\textbf{Nota para el docente.} Las tres formulaciones del triángulo son la síntesis del autor de la guía, no citas textuales de los pensadores. Si vas a citarlos en clase, remite a la obra original.
\end{infoband}

\milcestacion{milcVino}{Mi autoevaluación · marca una casilla por fila}

{\small
\setlength{\extrarowheight}{2.2pt}
\begin{tabularx}{\textwidth}{>{\RaggedRight\arraybackslash\bfseries}p{3.0cm}YYY}
\rowcolor{milcVino}\color{white}\bfseries Criterio & \color{white}\bfseries Lo logré & \color{white}\bfseries En proceso & \color{white}\bfseries Necesito apoyo\\[.6mm]
Comprensión del tema & \milcopcion{Explico las ideas con mis palabras.} & \milcopcion{Reconozco las ideas, falta precisión.} & \milcopcion{Necesito repasar lo básico.}\\[.8mm]
\rowcolor{milcGris}Pensamiento computacional & \milcopcion{Usé los pilares al armar mi producto.} & \milcopcion{Usé algunos pilares.} & \milcopcion{Necesito releer la estación 2.}\\[.8mm]
Aplicación y producto & \milcopcion{Mi producto cumple los criterios.} & \milcopcion{Está armado, con ajustes pendientes.} & \milcopcion{Aún está en construcción.}\\[.8mm]
\rowcolor{milcGris}Reflexión 5 dimensiones & \milcopcion{Respondí cada una con honestidad.} & \milcopcion{Respondí algunas con detalle.} & \milcopcion{Mis respuestas son muy generales.}\\[.8mm]
Triángulo de pensamiento & \milcopcion{Conecté las 3 voces con mi trabajo.} & \milcopcion{Conecté al menos una.} & \milcopcion{No las relaciono con mi proceso.}\\[.8mm]
\end{tabularx}}

\vspace{3mm}
\textbf{Hoy entendí que:}\\[1mm]
\linead{\linewidth}\\[1mm]
\linead{\linewidth}

\textbf{Mi compromiso para la próxima sesión es:}\\[1mm]
\linead{\linewidth}\\[1mm]
\linead{\linewidth}

\begin{softbox}{milcMostaza}{milcGris}{Cita de despedida · Marco Aurelio}
\textit{``El obstáculo en el camino se vuelve el camino. Lo que estorba al camino se vuelve camino.''}\\[1mm]
Cada error de hoy, bien observado, se convierte en pieza del camino que pisarás mañana.
\end{softbox}

\small
\begin{tabularx}{\textwidth}{>{\bfseries}p{3.6cm}Y}
\rowcolor{milcGradeDark!12}Nombre completo: & \linead{10cm}\\
Fecha: & \linead{4cm} \quad Curso/grupo: \linead{4cm}\\
Mi firma: & \linead{9cm}\\
\end{tabularx}
\normalsize

\begin{infoband}{milcGradeDark}
\textbf{Para la próxima sesión.} Esta semana, antes de enviar cada mensaje, pasa por los 3 filtros. \emph{¿Verdad? ¿Necesario? ¿Amable?}. La próxima clase es la \textbf{Sesión 10 --- Mi presentación digital}: vas a cosechar todo lo aprendido en una mini exposición de 2 minutos para tus compañeros. Hoy aprendiste a comunicar; la próxima clase muestras todo lo que aprendiste en P1 a tu audiencia.
\end{infoband}

\end{document}
